\documentclass[a4paper,12pt]{article}
\usepackage[T2A]{fontenc}
\usepackage[utf8]{inputenc}
\usepackage[russian]{babel}
\usepackage{amsmath,amssymb,mathtools}
\renewcommand{\familydefault}{\sfdefault}
\usepackage{icomma}
\usepackage[a4paper,margin=2.05cm]{geometry}
\usepackage{microtype}
\usepackage{tikz}
\usetikzlibrary{arrows.meta,positioning,shapes.geometric,calc}
\usepackage{tabularx,booktabs,longtable}
\usepackage{enumitem}
\usepackage{parskip}
\usepackage{fancyhdr}
\usepackage{setspace}
\usepackage{xcolor}
\usepackage[hidelinks]{hyperref}

\definecolor{accent}{HTML}{2E6F73}
\definecolor{darkaccent}{HTML}{1F4E52}
\definecolor{textgray}{HTML}{2F3437}
\definecolor{softgray}{HTML}{EEF2F2}

\setlength{\parindent}{0pt}
\setlength{\parskip}{0.45em}
\onehalfspacing
\setlist[itemize]{leftmargin=1.45em,itemsep=0.12em,topsep=0.15em}
\setlist[enumerate]{leftmargin=1.65em,itemsep=0.18em,topsep=0.15em}

\pagestyle{fancy}
\fancyhf{}
\renewcommand{\headrulewidth}{0pt}
\fancyfoot[C]{\small\color{textgray}06.10.26}

\newcommand{\sectiontitle}[1]{%
  \vspace{0.2em}%
  {\large\bfseries\color{darkaccent}#1}\par
  \vspace{0.08em}\hrule height 0.55pt\vspace{0.42em}%
}
\newcommand{\key}[1]{\textbf{\color{darkaccent}#1}}

\tikzset{
  flowstart/.style={ellipse, draw=accent, line width=0.9pt, fill=softgray, align=center, inner xsep=10pt, inner ysep=6pt},
  flowstep/.style={rounded corners=3pt, draw=accent, line width=0.9pt, fill=white, align=center, text width=0.72\textwidth, inner xsep=10pt, inner ysep=7pt},
  flowarrow/.style={-{Latex[length=3mm]}, line width=0.9pt, draw=accent}
}

\begin{document}

% ---------- TITLE ----------
\begin{titlepage}
\thispagestyle{empty}
\centering
\vspace*{2.5cm}

{\color{darkaccent}\bfseries\LARGE Чек-лист\par}
\vspace{0.55cm}
{\bfseries\Huge Приведение дробей\\к общему знаменателю\par}
\vspace{1.0cm}
{\Large Сложение, сравнение и упорядочивание дробей\par}

\vfill
{\large Ксения Васильченко\par}
\vspace{0.35cm}
{\large 06.10.26\par}
\vspace{1.1cm}
{\normalsize\color{textgray}Лёвин Артём Александрович эксклюзивно для Ксении Васильченко\par}

\vfill
\begin{tikzpicture}[remember picture,overlay]
  \draw[accent, line width=1.2pt]
    ([xshift=1.6cm,yshift=1.7cm]current page.south west)
      .. controls +(4.8cm,0.8cm) and +(-4.8cm,-0.8cm) ..
    ([xshift=-1.6cm,yshift=2.1cm]current page.south east);
\end{tikzpicture}
\end{titlepage}

% ---------- PAGE 2: CORE ----------
\sectiontitle{1. Главная идея}

Дроби удобно складывать, вычитать и сравнивать, когда у них \key{одинаковые знаменатели}. Для этого используем равные дроби:
\[
\frac{a}{b}=\frac{a\cdot k}{b\cdot k},\qquad k\ne0.
\]
То есть числитель и знаменатель умножаем на одно и то же число.

\textbf{Пример.}
\[
\frac13=\frac26,\qquad \frac12=\frac36,
\qquad
\frac13+\frac12=\frac26+\frac36=\frac56.
\]

\sectiontitle{2. Общий знаменатель через НОК}

Удобнее всего брать \key{НОК знаменателей}. Чтобы найти его через разложение на простые множители:
\begin{enumerate}
  \item разложите каждый знаменатель;
  \item выпишите множители одного числа;
  \item добавьте из остальных разложений только те множители, которых ещё не хватает;
  \item перемножьте их.
\end{enumerate}

\textbf{Пример: 24 и 32.}
\[
24=2\cdot2\cdot2\cdot3,
\qquad
32=2\cdot2\cdot2\cdot2\cdot2,
\]
\[
\text{НОК}(24,32)=2\cdot2\cdot2\cdot2\cdot2\cdot3=96.
\]

\sectiontitle{3. Дополнительный множитель}

После нахождения НОК для каждой дроби вычислите
\[
\text{дополнительный множитель}=\frac{\text{НОК}}{\text{знаменатель}}.
\]
Для знаменателей 24 и 32 это $96:24=4$ и $96:32=3$. Поэтому
\[
\frac1{24}=\frac4{96},\qquad
\frac1{32}=\frac3{96},
\qquad
\frac1{24}+\frac1{32}=\frac7{96}.
\]

\clearpage

% ---------- PAGE 3: OPERATIONS ----------
\sectiontitle{4. Сложение и вычитание дробей}

\textbf{Рабочий порядок:}
\begin{enumerate}
  \item найдите НОК знаменателей;
  \item найдите дополнительные множители;
  \item домножьте и числители, и знаменатели;
  \item выполните действие с числителями;
  \item сократите ответ, если это возможно.
\end{enumerate}

\textbf{Пример.}
\[
\frac1{18}+\frac1{24}.
\]
\[
18=2\cdot3\cdot3,\qquad 24=2\cdot2\cdot2\cdot3,
\qquad \text{НОК}(18,24)=72.
\]
\[
\frac1{18}+\frac1{24}
=\frac4{72}+\frac3{72}
=\frac7{72}.
\]

\sectiontitle{5. Если дробей несколько}

Тот же метод работает сразу для трёх и более дробей: ищем НОК \key{всех} знаменателей.

\textbf{Пример.}
\[
\frac1{12}+\frac1{15}+\frac1{18}.
\]
\[
\text{НОК}(12,15,18)=180,
\]
\[
\frac1{12}=\frac{15}{180},\qquad
\frac1{15}=\frac{12}{180},\qquad
\frac1{18}=\frac{10}{180},
\]
\[
\frac1{12}+\frac1{15}+\frac1{18}
=\frac{15+12+10}{180}
=\frac{37}{180}.
\]

\sectiontitle{6. Сокращение результата}

После вычислений задайте себе вопрос: \key{есть ли у числителя и знаменателя общий делитель?}
\[
\frac9{84}=\frac{9:3}{84:3}=\frac3{28}.
\]
Для быстрой проверки сначала удобно проверить делимость обоих чисел на $2$, затем на $3$, затем на $5$.

\clearpage

% ---------- PAGE 4: COMPARISON ----------
\sectiontitle{7. Сравнение дробей}

Если знаменатели одинаковы, больше та дробь, у которой больше числитель. При разных знаменателях сначала приведите дроби к общему знаменателю.

\textbf{Пример.}
\[
\frac5{12}\quad\text{и}\quad\frac7{18},
\qquad \text{НОК}(12,18)=36.
\]
\[
\frac5{12}=\frac{15}{36},\qquad
\frac7{18}=\frac{14}{36}.
\]
Так как $15>14$, то
\[
\boxed{\frac5{12}>\frac7{18}}.
\]

\sectiontitle{8. Расположение дробей по возрастанию}

Для нескольких дробей используйте один общий знаменатель для всех.

\textbf{Пример.}
\[
\frac9{16},\qquad \frac23,\qquad \frac{11}{21},
\qquad \text{НОК}(16,3,21)=336.
\]
\[
\frac9{16}=\frac{189}{336},\quad
\frac23=\frac{224}{336},\quad
\frac{11}{21}=\frac{176}{336}.
\]
Следовательно,
\[
\boxed{\frac{11}{21}<\frac9{16}<\frac23}.
\]

\sectiontitle{9. Типичные ошибки — короткая проверка}

\begin{tabularx}{\linewidth}{@{}>{\bfseries}p{0.33\linewidth}X@{}}
\toprule
Ошибка & Что проверить \\
\midrule
Неверное разложение & Перемножьте простые множители обратно: должно получиться исходное число. \\
Домножен только знаменатель & Дополнительный множитель умножает и числитель, и знаменатель. \\
Лишний множитель в НОК & Добавляйте только то, чего ещё не хватает. \\
Неверное сравнение & После общего знаменателя сравнивайте числители. \\
Несокращённый ответ & Проверьте общий делитель числителя и знаменателя. \\
\bottomrule
\end{tabularx}

\vspace{0.35em}
\textbf{Знаки:} $>$ — больше, $<$ — меньше. Знак $>$ можно мысленно достроить до буквы «В» как подсказку к слову «больше».

\clearpage

% ---------- FLOWCHART ----------
\thispagestyle{fancy}
\begin{center}
{\LARGE\bfseries\color{darkaccent}Алгоритм: общий знаменатель}\par
\vspace{0.9cm}

\begin{tikzpicture}[node distance=7mm]
  \node[flowstart] (start) {Начало};
  \node[flowstep, below=of start] (den) {Выпишите знаменатели всех дробей};
  \node[flowstep, below=of den] (factor) {Разложите знаменатели на простые множители};
  \node[flowstep, below=of factor] (lcm) {Найдите НОК знаменателей};
  \node[flowstep, below=of lcm] (mult) {Для каждой дроби найдите дополнительный множитель: $\text{НОК}:\text{знаменатель}$};
  \node[flowstep, below=of mult] (equiv) {Умножьте на него и числитель, и знаменатель};
  \node[flowstep, below=of equiv] (operate) {Выполните действие: сложите, вычтите или сравните числители};
  \node[flowstep, below=of operate] (reduce) {Если получили ответ-дробь, проверьте возможность сокращения};
  \node[flowstart, below=of reduce] (finish) {Готово};

  \draw[flowarrow] (start) -- (den);
  \draw[flowarrow] (den) -- (factor);
  \draw[flowarrow] (factor) -- (lcm);
  \draw[flowarrow] (lcm) -- (mult);
  \draw[flowarrow] (mult) -- (equiv);
  \draw[flowarrow] (equiv) -- (operate);
  \draw[flowarrow] (operate) -- (reduce);
  \draw[flowarrow] (reduce) -- (finish);
\end{tikzpicture}
\end{center}

\clearpage

% ---------- TRAINING ----------
\sectiontitle{Тренировочный блок — Путь А}

Решите 10 заданий. В вычислительных заданиях используйте общий знаменатель через НОК и сокращайте ответ, если это возможно.

\begin{enumerate}
  \item Приведите дроби $\dfrac38$ и $\dfrac5{12}$ к наименьшему общему знаменателю.

  \item Вычислите $\dfrac5{18}+\dfrac7{24}$.

  \item Вычислите $\dfrac{11}{20}-\dfrac7{30}$.

  \item Вычислите $\dfrac1{12}+\dfrac1{15}+\dfrac1{20}$.

  \item Сравните $\dfrac7{12}$ и $\dfrac58$.

  \item Расположите по возрастанию:
  \[
  \frac5{12},\qquad \frac7{18},\qquad \frac38.
  \]

  \item Определите, верно ли неравенство $\dfrac5{14}>\dfrac{13}{35}$.

  \item Вычислите $\dfrac7{24}+\dfrac5{36}$.

  \item Вычислите $\dfrac{13}{30}-\dfrac7{45}$.

  \item Вычислите $\dfrac5{12}+\dfrac7{18}-\dfrac19$.
\end{enumerate}

\clearpage

% ---------- ANSWERS ----------
\sectiontitle{Ответы}

\renewcommand{\arraystretch}{1.35}
\begin{center}
\begin{tabularx}{0.96\linewidth}{@{}cX@{}}
\toprule
\textbf{№} & \textbf{Ответ} \\
\midrule
1 & $\dfrac9{24}$ и $\dfrac{10}{24}$ \\
2 & $\dfrac{41}{72}$ \\
3 & $\dfrac{19}{60}$ \\
4 & $\dfrac15$ \\
5 & $\dfrac7{12}<\dfrac58$ \\
6 & $\dfrac38<\dfrac7{18}<\dfrac5{12}$ \\
7 & Нет, верно $\dfrac5{14}<\dfrac{13}{35}$ \\
8 & $\dfrac{31}{72}$ \\
9 & $\dfrac5{18}$ \\
10 & $\dfrac{25}{36}$ \\
\bottomrule
\end{tabularx}
\end{center}

\vfill
{\small\color{textgray}Лёвин Артём Александрович эксклюзивно для Ксении Васильченко}

\end{document}
